\section{\texorpdfstring{$R(1,3)$型包絡と楕円$E_6$量子化の同じ表示族}{R(1,3)型包絡と楕円E6量子化の同じ表示族}}
\label{sec:e6-r13-comparison}
本節の比較対象は、中心一次元を持つ\emph{次数付き包絡と中心三次直線の対}
$(B,kG)$である。任意周期付き$\mathbb Z$-代数の全同型や、
構造対象を忘れた圏の同値を比較しているのではない。
Pym \cite[\S3.4, (9)]{PymR13}の$R(1,3)$型は三次元Sklyanin環の中心一次拡大であり、
その包絡の一般商は七次元である。
中心三次商まで選んだ八次元の族と、本稿の楕円$E_6$型を別々の成分候補として数えてはならない。
以下では次元だけでなく関係式を比較する。

\subsection{巡回ポテンシャルの交代部分とJacobian表示}
$V=kx\oplus ky\oplus kz$とし、巡回語は自由代数の交換子部分空間による商で考える。
$\operatorname{sym}$は通常の対称化を巡回語空間へ写す写像とする。
例えば$\operatorname{sym}(xyz)=(xyz+xzy)/2$である。
\begin{lemma}[交代係数が非零の巡回chart]
\label{lem:e6-r13-cyclic-chart}
定数項を除いた次数三以下の巡回ポテンシャルは一意に
\[
 \Phi=\alpha(xyz-xzy)+\operatorname{sym}(p_3+p_2+p_1),
 \qquad p_i\in\operatorname{Sym}^iV
\]
と書ける。$\alpha\ne0$なら、$f=-\alpha^{-1}(p_3+p_2+p_1)$と置いて
\[
 \alpha^{-1}\Phi=xyz-xzy-\operatorname{sym}(f)
\]
と規格化できる。そのJacobian代数のRees包絡は
\begin{equation}\label{eq:e6-r13-envelope}
 [t,x]=[t,y]=[t,z]=0,\qquad
 [y,z]=\operatorname{Sym}(\partial_xF_f),\quad\text{および巡回式},
\end{equation}
ただし$F_f=p'_3+t p'_2+t^2p'_1$、$f=p'_3+p'_2+p'_1$であり、
右辺の$\operatorname{Sym}$は\emph{自由代数内の}対称化である。
従ってこれは$R(1,3)$中心拡大のJacobian表示そのものである。
\end{lemma}
\begin{proof}
三次巡回語は十一個あり、十個の対称三次単項式と
$xyz-xzy$が基底をなす。実際、三文字が全て異なる場合の二軌道は
$xyz,xzy$であり、対称和と差へ分解できる。
それ以外は各多重次数に一軌道しかない。二次・一次巡回語は対称語である。
巡回微分について
$\partial_x^{\rm cyc}(xyz-xzy)=yz-zy$および
$\partial_x^{\rm cyc}\operatorname{sym}(p)=\operatorname{Sym}(\partial_xp)$
が単項式ごとに成り立つ。よって三関係は
$[y,z]-\operatorname{Sym}(\partial_xf)=0$とその巡回式である。
$t$による斉次化で式\eqref{eq:e6-r13-envelope}を得る。
ポテンシャルの非零スカラー倍は関係イデアルを変えないので、
この比較は包絡だけでなく、同じ元$G$を選んだ中心三次商も保つ。
\end{proof}

この恒等式自体はAS正則性を含意しない。
以下では本稿の全次数平坦なAS4表示開集合と、
既知の中心拡大の良い範囲 \cite{LBS,EG,PymR13}に制限する。
また$f_3=0$等の退化を、滑らかな最高三次を持つ一般chartに含めない。

\subsection{八変数族における明示的な係数変換}
定理\ref{thm:eight-formal-central-lift}の記号を使い、
$p=1+\kappa\lambda$とする。本稿のポテンシャルは
\[
 \Phi=pxyz-q_0xzy+\lambda(x^3+\eta^3y^3+\zeta^3z^3)
       -\tfrac12(x^2+y^2+z^2)-ax-by-cz.
\]
補題\ref{lem:e6-r13-cyclic-chart}を適用すると
\begin{align}
 \alpha&=(p+q_0)/2,\notag\\
 f&=-\frac{2}{p+q_0}\left(
 (p-q_0)xyz+\lambda(x^3+\eta^3y^3+\zeta^3z^3)
 -\tfrac12(x^2+y^2+z^2)-ax-by-cz\right).
 \label{eq:e6-r13-f}
\end{align}
$\lambda=0$で$p+q_0=1+q_0\ne0$なので、この規格化は特殊点でも正則である。

\begin{proposition}[一般包絡の$R(1,3)$表示]
\label{prop:e6-r13-generic-envelope}
命題\ref{prop:algebraic-presentation-flat}の底を、基準点を含む開部分へ縮める。
その上の稠密開部分$\lambda\ne0$では、
式\eqref{eq:e6-r13-f}の最高斉次三次$f_3$は滑らかであり、
本稿の八変数族は$R(1,3)$型包絡と中心三次直線の対の族へ入る。
$\lambda=0$では$f_3$は非零スカラー倍の$xyz$となる。
これは\emph{切断の反標準楕円境界が特異になる}という意味ではない。
\end{proposition}
\begin{proof}
$A x^3+B y^3+C z^3+Lxyz$で$ABC\ne0$の場合、
三偏微分に共通射影零点があることは$L^3+27ABC=0$と同値である。
座標の一つが零なら偏微分から全座標が零となる。
全て非零なら三偏微分式を掛けて必要条件を得る。
逆は代数閉性と座標スカラー変換によりHesse三次の特異点から従う。
本稿では
\[
 A=\lambda,\quad B=\lambda\eta^3,\quad C=\lambda\zeta^3,
 \quad L=p-q_0,
 \qquad \Delta_f=(p-q_0)^3+27\lambda^3\eta^3\zeta^3.
\]
$\eta\zeta\ne0$かつ$\Delta_f|_{\lambda=0}=(1-q_0)^3\ne0$であるから、
$p+q_0$と$\Delta_f$をさらに可逆化すれば主張の稠密開部分を得る。
包絡の同定は補題による恒等式であり、中心元$G$も同じ環内の同じ元として保たれる。
全次数flat性・AS4性・中心性は既証の命題
\ref{prop:algebraic-presentation-flat}を用いる。
$\lambda=0$においても、切断によって得る楕円曲線は
第\ref{sec:seven}節の$E_u$であり、$f_3=0$とは別の曲線である。
\end{proof}

\subsection{包絡の七次元と中心切断を含む八次元}
最高三次が滑らかな$f=f_3+f_2+f_1$の係数空間は十九次元である。
中心一次直線を保つgraded同型を共通次数スカラーで規格化すると、
生成元の変換は$\mathbf x\mapsto A\mathbf x+\mathbf b t$となり、
\[
 f\longmapsto\frac{f(A\mathbf x+\mathbf b)-f(\mathbf b)}{\det A}
\]
という$\operatorname{Aff}_3$作用を得る。
この同型群oidと一般安定化群の有限性は
\cite[定理 \texttt{thm:r13-affine}]{CompanionAmbient}で証明されている。
従って包絡の一般商の次元は$19-12=7$である。

中心三次については、先頭AS3環$S_0$が
$Z(S_0)_1=Z(S_0)_2=0$, $Z(S_0)_3=k\omega$を持ち、
$\omega$が中心三次元$\Omega_B$へ持ち上がる良い範囲を使う。
最高active記号を順に引くことにより
\[
 Z(B)_3=k\Omega_B\oplus kt^3.
\]
よって三次直線の選択は$\PP^1$であり、$[t^3]$を除けば
$G=\Omega_B+D t^3$と書ける。
有限な一般射影安定化群で割ってもこの一次元は残るので、
中心三次商の一般表示商は$7+1=8$次元である。
これが\cite[命題 \texttt{prop:center}]{CompanionAmbient}の一次元ファイバーであり、
本稿の独立な中心値$D$に対応する。

\begin{corollary}[同じ表示成分の一般部分と三捩れ特殊部分]
\label{cor:e6-r13-same-presentation}
交代係数非零、滑らかな最高三次、上記の低次数中心条件と有限安定化群を満たす
共通開部分では、$R(1,3)$型包絡の中心三次商と本稿の一般楕円$E_6$量子化は
同じ表示incidenceに属する。
定理\ref{thm:formal-potential-incidence-smooth}の基準点を通る唯一の局所既約成分は、
この一般表示族の閉包である。
従って本稿の八次元表示成分を、包絡の$R(1,3)$成分と別の曲面成分として数えない。
\end{corollary}
\begin{proof}
交代係数を1へ規格化する操作は、巡回potential modulo scalarの開部分と
十九係数の$f$の空間を同定する。補題の関係イデアルは同じなので、
中心元の線型方程式も同じである。低次数中心が一定rank二の開部分では、
中心直線のincidenceはこの既約な係数開部分上の$\PP^1$束となり既約である。
命題\ref{prop:e6-r13-generic-envelope}と既証の一般中心条件により、
八変数族の幾何学的一般点はこのincidenceへ入る。
一方、指定点の全表示incidenceの完備局所環は
定理\ref{thm:formal-potential-incidence-smooth}により二十二変数形式冪級数環で、
座標・二スカラーの十四次元を除いた局所表示次元は八である。
同定理で八変数族はこの局所表示を尽くすため、
その稠密開部分の閉包が唯一の局所成分を与える。
\end{proof}

\begin{remark}[三つの異なる比較問題]\label{rem:e6-r13-scope}
包絡の七次元、固定三捩れの\emph{曲面}の七次元、
三捩れを動かす曲面の八次元を区別する。
固定三捩れでは包絡は$(a,b,c)$の三次元、中心直線は$\PP^4$の四次元であり、
一般包絡上の$7+1$という分解とは異なる。
また$R(1,3)$という名称の一致は、AOUのAS曲面であることを含意しない。
第\ref{sec:as-surface-scope}節のpoint-scheme障害は固定三捩れの族に限られ、
横断した八次元族全体の非AS性には使えない。
最後に、本節は次数付き表示の比較である。
全intrinsic H3スタックの開性・相対再構成・降下・全境界の同定については、
定理\ref{thm:conditional-component}の追加仮定を維持する。
\end{remark}

補題の巡回基底、次数一から三の全十九単項式に対する五十七個の微分恒等式、
式\eqref{eq:e6-r13-f}の全係数変換、Hesse判別条件の必要恒等式と特殊点の単元性を
\texttt{computations/verify\_r13\_comparison.py}で厳密検算した。
これは有限標本からの外挿ではないが、AS4性・H3性や圏の相対降下を検算だけで証明するものでもない。
